\documentclass[12pt,a4paper]{article}
\usepackage[T2A,T1]{fontenc}
\usepackage[utf8]{inputenc}
\usepackage[english,russian]{babel}
\usepackage{XCharter}
\usepackage{amsmath,amssymb,mathtools}
\usepackage{array,longtable,booktabs}
\usepackage[margin=24mm]{geometry}
\usepackage[hidelinks]{hyperref}
\setlength{\parindent}{0pt}
\setlength{\parskip}{0.55em}
\setlength{\LTleft}{0pt}
\emergencystretch=2em
\setlength{\LTright}{0pt}
\newcolumntype{P}[1]{>{\raggedright\arraybackslash}p{#1}}
\title{От посылки к заключению}
\author{}
\date{}
\begin{document}
\maketitle
\begin{center}\emph{Конспект о доказательстве, необходимом следовании и споре Н. П. Брусенцова с классической логикой}\end{center}

\section{Пример: от посылок к заключению}

Из двух посылок «всякое целое число, которое делится на 4, чётно» и «12 делится на 4» получаем заключение «12 чётно». Здесь возникают два вопроса. По какому правилу мы делаем этот шаг? Чем обоснована сама общая посылка о делимости? Первый вопрос касается правил вывода. Второй требует обратиться к свойствам целых чисел.

\section{Что участвует в рассуждении}

\begin{longtable}{P{.23\linewidth}P{.70\linewidth}}
\toprule
\textbf{Понятие} & \textbf{Значение в этом конспекте} \\
\midrule
Признак & Свойство предмета, например чётность числа. Предикат «… чётно» выражает этот признак. Подставив число 12, получаем утверждение «12 чётно». \\
Высказывание & То, что утверждается и при заданном истолковании может быть истинно или ложно. Буквы P и Q ниже обозначают целые высказывания. \\
Посылка и заключение & Посылка — высказывание, принятое в качестве основания рассуждения. Заключение — высказывание, которое из этих оснований выводят. В следующем рассуждении уже доказанное заключение может стать посылкой. \\
Правило вывода & Схема допустимого перехода от уже принятых положений к следующему положению. \\
Доказательство & Цепочка обоснованных шагов от допущений, аксиом или ранее доказанных положений к тезису. \\
Модель & В формальной логике — математическая интерпретация языка, в которой можно проверить истинность утверждений. В логике высказываний достаточно назначить исходным буквам значения «истина» и «ложь». \\
\bottomrule
\end{longtable}

\begin{description}
\item[\(P\to Q\)] Условная формула («если \(P\), то \(Q\)»). Её классические условия истинности показывает таблица в следующем разделе.
\item[\(P\models Q\)] Семантическое следование: во всякой модели, где \(P\) истинно, \(Q\) тоже истинно.
\item[\(P\vdash Q\)] Выводимость: в выбранном исчислении существует доказательство \(Q\) из посылки \(P\).
\end{description}
Это три разные записи. Для стандартного классического исчисления высказываний, если семантику и правила вывода задать согласованно, семантическое следование и выводимость совпадают. Первое устанавливают отсутствием контрпримера, второе — существованием доказательства. Это говорит о формальной корректности перехода при данных посылках, но не устанавливает истинность самих посылок. \href{https://plato.stanford.edu/entries/logical-consequence/}{Stanford Encyclopedia of Philosophy, Logical Consequence}

Сам переход к заключению не является формулой \(P\to Q\). Заключение — это само высказывание \(Q\); правило вывода указывает, когда его разрешено записать из посылок. Например, из \(P\to Q\) и \(P\) по \emph{modus ponens} выводят \(Q\). Формула \(P\to Q\) может входить в доказательство, но не заменяет ни правило, ни заключение.

\subsection{Порядок проверки рассуждения}

Перед оценкой аргумента полезно последовательно задать несколько разных вопросов. Это рабочий порядок анализа, а не ещё одно правило вывода.

\begin{enumerate}
\item \textbf{Что утверждается и на каком основании?} Выпишите тезис и каждую посылку отдельно; восстановите пропущенные посылки и отметьте, какие из них сами требуют обоснования.
\item \textbf{О каких предметах и в каком смысле говорится?} Зафиксируйте область рассмотрения, значение терминов и слова «все», «некоторые», «ни один». Сохраняйте направление: «все \(X\) — \(Y\)» и «все \(Y\) — \(X\)» — разные утверждения. Не переносите свойство группы на каждого её участника.
\item \textbf{Каков вид перехода?} Различите категорический силлогизм, вывод по условной посылке, обобщение по наблюдаемым случаям и объяснение причины. Восемь традиционных правил применяют к простому категорическому силлогизму с двумя посылками и тремя терминами. Для \emph{modus ponens} нужны две посылки: \(P\to Q\) и \(P\). Если из общих посылок выводится заключение со словом «некоторые», в традиционной силлогистике отдельно проверьте предположение о существовании предметов нужного класса.
\item \textbf{Следует ли заключение?} Для дедуктивного вывода укажите формальную систему, её язык и правила. Одна допустимая модель, в которой все посылки истинны, а заключение ложно, опровергает семантическое следование; чтобы подтвердить его, надо показать, что такой модели нет. Выводимость устанавливают построением доказательства по правилам выбранного исчисления. Если исчисление корректно и полно относительно заданной семантики, семантическое следование и выводимость совпадают. Таблица истинности проверяет формулы классической логики высказываний, но при проверке силлогизма надо сохранить кванторы, предикаты и средний термин.

Для обобщения по наблюдениям оценивают охват и качество случаев; для причинного объяснения — данные и возможные альтернативные причины.
\item \textbf{Обоснованы ли посылки?} Это отдельная проверка: правильная форма сама по себе не делает посылки истинными. В математике укажите, какие положения приняты как аксиомы или допущения, а какие уже доказаны; в эмпирическом рассуждении проверьте данные и границы их применимости.
\item \textbf{Не требуется ли дополнительная связь?} Формальное следование, истинность посылок и причинное объяснение — разные вопросы. Содержательное следование Н. П. Брусенцова и релевантное следование в современной логике также нельзя считать одним понятием без сопоставления определений. Таблица истинности сама по себе не устанавливает причинность. При разборе теории Н. П. Брусенцова укажите конкретную публикацию, определения, универсум и постулаты; классические законы и обозначения нельзя переносить на неё без обоснования. Если называете логическую ошибку, укажите шаг рассуждения и нарушенное условие: само название не заменяет проверку.
\end{enumerate}

\section{Что показывает таблица истинности}

Условную формулу \(P\to Q\) читают «если P, то Q». P называют \emph{антецедентом}, Q — \emph{консеквентом}. Первые два столбца показывают значения двух исходных высказываний. Последний показывает истинность всего условного высказывания \(P\to Q\) при данном сочетании значений. Одна эта таблица не кодирует причинную, временную или объяснительную связь между P и Q.

\begin{longtable}{P{.18\linewidth}P{.43\linewidth}P{.27\linewidth}}
\toprule
\textbf{P} & \textbf{Q} & \textbf{\(P\to Q\)} \\
\midrule
Истина & Истина & Истина \\
Истина & Ложь & \textbf{Ложь} \\
Ложь & Истина & Истина \\
Ложь & Ложь & Истина \\
\bottomrule
\end{longtable}

Например, P означает «лампа включена», Q — «дверь открыта». Если лампа выключена, формула \(P\to Q\) истинна при любом положении двери. Но есть и возможная ситуация, когда лампа включена, а дверь закрыта. Это контрпример к выводу «лампа включена, следовательно, дверь открыта». Одна ложная посылка означает, что конкретная оценка не подтверждает рассуждение; противоречивый набор посылок означает, что нет ни одной оценки, где все они истинны одновременно. Поэтому в классической логике из противоречивого набора формально следует любое заключение: не нашлось модели, которая удовлетворяет посылкам и опровергает заключение.

\textbf{Критерий следования в классической логике.} Из единственной посылки \(P\) следует \(Q\), если нет допустимой оценки, где \(P\) истинно, а \(Q\) ложно; для классической логики это равносильно тому, что \(P\to Q\) истинно при всех оценках. Для нескольких посылок \(\Gamma\) (здесь \(\Gamma\) — множество посылок) проверяют все оценки, в которых истинны сразу все посылки: ни в одной из них заключение \(Q\) не должно быть ложным. \href{https://plato.stanford.edu/entries/logic-propositional/}{Stanford Encyclopedia of Philosophy, Propositional Logic}

\section{Как математик делает шаг доказательства}

\textbf{Введение импликации.} Возьмём произвольное целое \(n\). Временно предположим, что \(n\) делится на 4. Тогда по определению \(n=4k\) для некоторого целого \(k\). Перепишем равенство в виде \(n=2(2k)\). Множитель \(2k\) целый, а значит, \(n\) чётно. Теперь снимаем временное предположение и получаем условное утверждение «если \(n\) делится на 4, то \(n\) чётно». Выбор \(n\) был произвольным, поэтому утверждение верно для каждого целого числа.

\textbf{Применение импликации (modus ponens).} Подставим 12 вместо \(n\) в доказанную общую теорему. Получим первую посылку «если 12 делится на 4, то 12 чётно». Равенство \(12=4\cdot3\) обосновывает вторую посылку «12 делится на 4». Теперь по правилу \emph{modus ponens} записываем заключение «12 чётно». В записи \textbf{\(P\to Q,\ P\vdash Q\)} знак \(\vdash\) сообщает, что заключение выводимо из указанных посылок. Название правила указывает, как именно сделан этот шаг. \href{https://builds.openlogicproject.org/content/first-order-logic/natural-deduction/propositional-rules.pdf}{Open Logic Project, правила для импликации}

\textbf{Необходимое и достаточное условия.} Делимости на 4 \emph{достаточно} для чётности. Чётность \emph{необходима} для делимости на 4. Оба предложения выражают доказанное направление «делится на 4 \(\to\) чётное». Обратное утверждение звучит иначе: «всякое чётное число делится на 4». Оно ложно, поскольку 6 чётно и не делится на 4. Чтобы назвать условие необходимым и достаточным, требуется доказать оба направления.

\textbf{Преобразование выражения и преобразование условия.} Запись \(4k=2(2k)\) позволяет заменить одно выражение другим с тем же значением. При решении уравнения сравниваются уже условия на неизвестное. Над действительными числами \(x-1=2\) равносильно \(x=3\), потому что прибавление единицы к обеим частям обратимо. Из \(x=3\) получается \(x^2=9\), но последнее уравнение допускает также \(-3\). После такого одностороннего шага подстановка в исходное уравнение отсеивает лишний корень.

\section{Что у Аристотеля означает «следовательно»}

В «Первой аналитике» (I.1, 24b18–22) Аристотель определяет \emph{силлогизм}. По смыслу его определения, это рассуждение, в котором из положенных посылок с необходимостью получается нечто отличное от них именно благодаря тому, что они положены. В обычном категорическом силлогизме связь можно проследить через \textbf{средний термин}, общий двум посылкам. \href{https://classics.mit.edu/Aristotle/prior.1.i.html}{Аристотель, Первая аналитика, I.1 и I.4}

\begin{quote}
Все люди смертны.\par
Все греки — люди.\par
Следовательно, все греки смертны.\par
\end{quote}

Средний термин «люди» связывает обе посылки. Возьмём произвольного грека. По второй посылке он человек, по первой — смертен. Поэтому заключение относится к каждому греку.

Если заменить три предложения тремя независимыми буквами \(P\), \(Q\), \(R\), получится лишь грубая пропозициональная схема \((P\land Q)\to R\): знак \(\land\) означает «и», а \(\to\) — «если…, то…». Схема ложна при \(P=Q=\text{истина}\) и \(R=\text{ложь}\). Это не контрпример к силлогизму. Такая замена стирает внутреннее устройство категорических посылок: квантор «все», предикаты и общий средний термин. Полная современная запись сохраняет их:
\[
\forall z\,(G(z)\to H(z)),\quad \forall z\,(H(z)\to M(z))
\;\models\; \forall z\,(G(z)\to M(z)),
\]
где \(\forall\) читается «для каждого предмета», \(G\) — «грек», \(H\) — «человек», \(M\) — «смертен», а \(z\) пробегает предметы области. При любой интерпретации, удовлетворяющей обеим посылкам, всякий грек является человеком и потому смертен. Поэтому неверно говорить, что исходные значения \(P,Q,R\) просто «исключаются словом все»: это буквенная абстракция теряет ту структуру, которая и обеспечивает вывод. \href{https://plato.stanford.edu/entries/aristotle-logic/}{Stanford Encyclopedia of Philosophy, логика Аристотеля}

Во «Второй аналитике» Аристотель задаёт более сильный вопрос. Можно знать, \textbf{что} произошло лунное затмение. Чтобы знать, \textbf{почему} оно произошло, нужно установить, что Земля преградила солнечному свету путь к Луне. В доказательстве причины средний термин выражает объясняющее основание. Здесь появляются два разных «почему». На вопрос «почему мы вправе принять заключение?» отвечает связь заключения с посылками. На вопрос «почему происходит затмение?» отвечает причина самого явления. Указание этой причины предъявляет дополнительное требование к посылкам объясняющего доказательства. \href{https://classics.mit.edu/Aristotle/posterior.1.i.html}{Вторая аналитика, I.2 и I.13} · \href{https://classics.mit.edu/Aristotle/posterior.2.ii.html}{II.2 и II.8, пример затмения}

\section{Силлогизм и математическая логика: что меняется у Буля и Кэрролла}

\subsection{Понятия, нужные для сравнения}

\begin{description}
\item[Класс.] Совокупность предметов, к которым применим некоторый термин. Например, класс греков — все предметы, о которых верно «это грек».
\item[Категорический силлогизм.] Вид силлогизма с посылками и заключением о связи между терминами («все», «ни один», «некоторые»). В обычной учебной схеме есть две посылки и заключение; средний термин входит в обе посылки и связывает крайние термины.
\item[Математическая логика.] Изучение формальных языков, моделей и правил доказательства математическими средствами. Это широкая область с разными способами записи и проверки вывода, а не одна отдельная связка или исчисление.
\item[Алгебра логики Буля.] Способ формально рассуждать о классах и отношениях между ними с помощью символов и алгебраических преобразований. Буль в книге 1847 года сначала разбирал силлогистику, а затем развивал методы для более широких задач; в \href{https://www.gutenberg.org/ebooks/15114}{книге 1854 года} он обобщил анализ на рассуждения со многими посылками и классами. Поэтому его подход не устраняет силлогизмы, а даёт им алгебраическую форму и расширяет круг задач. \href{https://plato.stanford.edu/entries/algebra-logic-tradition/}{Stanford Encyclopedia of Philosophy, The Algebra of Logic Tradition}; \href{https://www.gutenberg.org/ebooks/36884}{George Boole, The Mathematical Analysis of Logic (1847)}
\item[Классическая логика высказываний.] Формальная система, где целые высказывания обозначают буквами \(P,Q,\ldots\), а связки оценивают по таблицам истинности. Такая запись может скрыть устройство исходных высказываний — например, квантор «все» и общий средний термин силлогизма.
\item[Материальная импликация.] В логике высказываний \(P\to Q\) — классическая истинностная связка для целых высказываний: она ложна только при истинном \(P\) и ложном \(Q\); таблица приведена в разделе 3. В алгебре классов Н. П. Брусенцова \(x\to y\) — соответствующее двузначное классовое выражение. Эти записи имеют разные типы объектов, хотя используют аналогичную связку. Ни одна сама по себе не кодирует смысловую, причинную или объяснительную связь. Историческую алгебру классов Буля также нельзя без оговорки отождествлять с современной таблицей для \(P\to Q\).
\item[Правило вывода.] Разрешённый переход от посылок к заключению. Например, правило \emph{modus ponens} позволяет из \(P\to Q\) и \(P\) вывести \(Q\). Само правило — не ещё одно условное высказывание.
\end{description}

\subsection{В чём состоит переход, а в чём — не разрыв}

В силлогизме важна внутренняя структура категорических высказываний: «все», «некоторые», термины и их распределение по посылкам. При грубой замене целых посылок буквами \(P,Q,R\) и записи \((P\land Q)\to R\) эта структура пропадает. В записи через предикаты и кванторы её можно сохранить — как в разборе примера с греками выше. Значит, здесь есть разница между уровнями формализации, но не противопоставление «силлогизмы или математическая логика».

Исторически Буль не открыл разрыв с силлогизмами: он алгебраизировал традиционные формы рассуждения и расширил их средствами для многих посылок и классов. Кэрролл тоже не отказался от силлогистики. Его книга \emph{Symbolic Logic, Part I: Elementary} (1896) посвящена элементарной логике и, в частности, диаграммному решению категорических силлогизмов. \href{https://iep.utm.edu/lewis-carroll-logic/}{Internet Encyclopedia of Philosophy, Lewis Carroll: Logic}; \href{https://archive.org/details/symboliclogicpa00carrgoog}{Lewis Carroll, Symbolic Logic (1896)}

Особую проблему Кэрролл показывает в диалоге «Что Черепаха сказала Ахиллу» (1895). Черепаха принимает посылки \(A\) и \(B\), а затем и условное утверждение «если \(A\) и \(B\), то \(Z\)», но всё ещё не делает вывод \(Z\). Если Ахилл добавляет ещё одно условное утверждение, Черепаха требует следующего. Так возникает регресс: правило, применённое как шаг вывода, не заменяется добавлением того же содержания в список посылок. Это не опровержение \emph{modus ponens}; диалог подчёркивает различие между условной формулой и правилом её применения. \href{https://fair-use.org/mind/1895/04/what-the-tortoise-said-to-achilles}{Lewis Carroll, “What the Tortoise Said to Achilles” (1895)}

\subsection{Два разных смысла слова «парадокс»}

\textbf{Парадоксы материальной импликации} — это классически верные схемы, которые кажутся странными, если читать \(P\to Q\) как обещание содержательной связи:
\[
\neg P \models (P\to Q), \qquad Q \models (P\to Q).
\]
Первая схема верна потому, что материальная импликация истинна при ложном \(P\); вторая — потому, что она истинна при истинном \(Q\). Например, «если лампа включена, то дверь открыта» истинно по этой таблице, когда лампа выключена, независимо от положения двери. Это не формальное противоречие: парадоксально здесь расхождение таблицы истинности с привычным смыслом естественного «если». \href{https://plato.stanford.edu/entries/logic-conditionals/}{Stanford Encyclopedia of Philosophy, The Logic of Conditionals}

\textbf{Парадокс Кэрролла} касается другого: можно ли свести сам переход от посылок к заключению к ещё одной посылке-условию. Он показывает, почему формула \(P\to Q\), семантическое следование \(P\models Q\) и вывод по правилу \(P\to Q,\ P\vdash Q\) отвечают на разные вопросы. В следующем разделе это различение понадобится при чтении записи Н. П. Брусенцова \(x\Rightarrow y\): она обозначает предложенное отношение между терминами, а не автоматически материальную импликацию или правило вывода.

\section{Позиция Н. П. Брусенцова: материальная импликация и необходимое следование}

\subsection{Тезис автора и его оговорка}

Н. П. Брусенцов не оставляет открытым вопрос, выражает ли обычная двухзначная материальная импликация нужное ему отношение. В статье 2010 года он пишет, что необходимое следование «выродилось в „материальную импликацию“»; на следующей странице уточняет, что материальная импликация «не тождественна аристотелеву необходимому следованию» и «равнозначна ему только в подчинённом ПСП универсуме Аристотеля УА» (PDF с. 1–2, печатные с. 88–89). Следовательно, его тезис — не безусловное отсутствие всякой связи между импликацией и следованием: в принятом универсуме УА автор утверждает их равнозначность, а вне этой предпосылки отрицает тождество. \href{https://ternarycomp.cs.msu.ru/Papers/brusentsov_sborn2010_2.pdf}{«Воссоздание аристотелевой безукоризненной силлогистики»}

Содержательность он объясняет так: «Качество \(y\), присущее \(x\)-вещам, необходимо содержится в их сущности»; для противоположных \(x'\)-вещей оно «несущественно (может быть, а может не быть)» (2008, PDF с. 1). Это точная формулировка авторского философского основания, а не независимо доказанный факт о причинной или физической связи. В той же работе он называет материальную импликацию «не отношением, а типичной химерой» (PDF с. 4; соответствие этой страницы конкретной печатной странице не установлено). Эта оценка принадлежит автору и не является общепринятым результатом. \href{https://ternarycomp.cs.msu.ru/Papers/Alrek.pdf}{«Алгебраическая реконструкция силлогистики»}

Формальная часть возражения касается статусов сочетаний терминов. Если в нашем контрольном расчёте подставить \(x\to y\equiv x'\lor y\) в обычную двузначную оценку, получаем \(+-++\) в порядке \(xy,xy',x'y,x'y'\): ложна только ячейка \(xy'\). Это наше вычисление по классической таблице, а не профиль, который Н. П. Брусенцов вводит в источнике. В главе 2005 года он отдельно представляет необходимое следование нормальной формой \(V_{xy}V'_{xy'}V_{x'y'}\), а вектором \(+-0+\) кодирует статусы четырёх сочетаний: необходимое включение, исключение, привходящий статус и включение. В статье 2006 года он называет форму \(xy\lor\sigma x'y\lor x'y'\) «аналогом материальной импликации с устранёнными парадоксами»; сам текст 2006 года сжимает объяснение до этой формулировки. Подробное основание для привходящего статуса \(x'y\) дано в главе 2005 года, поэтому его не следует приписывать статье 2006 года как её собственное развернутое доказательство. \href{https://ternarycomp.cs.msu.ru/Papers/Rea.pdf}{«Реанимация аристотелевой силлогистики» (PDF с. 1–3)}; \href{https://ternarycomp.cs.msu.ru/Papers/Soder_sled.pdf}{«Содержательное следование, сосуществование противоположностей и формальные импликации» (PDF с. 1)}

Это не означает, что автор отрицает всякую связь между импликацией и своим отношением. В работе 2010 года он оговаривает: материальная импликация и необходимое следование равнозначны в универсуме \(\text{УА}\), подчинённом принципу сосуществования противоположностей. Эквивалентность, таким образом, зависит от дополнительного постулата; она не означает, что обычная материальная импликация без этой предпосылки уже выражает содержательное следование. \href{https://ternarycomp.cs.msu.ru/Papers/brusentsov_sborn2010_2.pdf}{«Воссоздание аристотелевой безукоризненной силлогистики» (печатные с. 88–89; PDF с. 1–2)}

\subsection{Не смешивать материальную и строгую импликации}

В главе 2005 года конкретный разбор начинается не с материальной, а со \textbf{строгой импликации} Льюиса. Автор специально оговаривает: «несуществование \(xy'\)-вещей не означает еще следования \(y\) из \(x\)» (PDF с. 1, печатная с. 140). Причина — пустота \(xy'\) совместима с отсутствием самих \(x\)-вещей или \(y'\)-вещей. Поэтому он добавляет условия существования \(x\) и \(y'\), записывает отношение как \(V_xV'_{xy'}V_{y'}\), а совершенную нормальную форму — как \(V_{xy}V'_{xy'}V_{x'y'}\). На PDF с. 2 он уточняет, что сочетание \(x'y\) не включено и не исключено с необходимостью, а имеет привходящий статус. Это отдельный аргумент о строгой импликации Льюиса, а не определение материальной импликации. \href{https://ternarycomp.cs.msu.ru/Papers/Rea.pdf}{«Реанимация аристотелевой силлогистики»}

В том же изложении \(V_x\) и \(V_{y'}\) выражают существование вещей, которым присущи соответствующие термины; \(V'_{xy'}\) — несуществование вещей с сочетанием \(xy'\). Это операции авторской модели, а не обычные квантор и отрицание из логики высказываний.

\newpage
\subsection{Термины, отношение и его представление}

В примере \(x\) — «носит очки», \(y\) — «студент среди участников собрания». Штрих в \(x'\) и \(y'\) обозначает противоположный термин. В этом примере его удобно понимать как отсутствие признака («без очков», «не студент»), но в общей модели противоположный термин не обязательно совпадает с обычным логическим отрицанием. Записи \(xy,xy',x'y,x'y'\) обозначают четыре группы предметов с соответствующими сочетаниями признаков.

\(x\Rightarrow y\) у Н. П. Брусенцова — отношение необходимого содержательного следования \(y\) из \(x\) между терминами. Его нельзя определять как сам профиль \(+-0+\): профиль в главе 2005 года служит алгебраическим представлением этого отношения через статусы четырёх сочетаний вещей. Это различие между отношением и его формальной записью.

Нужно также различать типы выражений. \(P\to Q\) — материальная импликация между целыми высказываниями; \(x\to y\) — её двузначный аналог в алгебре терминов и классов; \(x\Rightarrow y\) — отношение, которое Н. П. Брусенцов называет необходимым содержательным следованием. \(\Gamma\models Q\) сообщает о сохранении истинности во всех моделях, а \(\Gamma\vdash Q\) — о наличии вывода в выбранном исчислении. Это разные объекты и вопросы; ни одна запись не является синонимом другой.

\subsection{Профиль \(+-0+\) в примере с собранием}

В главе «Реанимация аристотелевой силлогистики» (2005) Н. П. Брусенцов представляет необходимое следование профилем сочетаний \(xy,xy',x'y,x'y'\) именно в таком порядке. Плюс кодирует необходимое включение подкласса, минус — необходимое исключение, а ноль в третьей позиции — привходящий статус: сочетание \(x'y\) может быть включено или исключено в зависимости от обстоятельств. В тексте автор прямо связывает нормальную форму \(V_{xy}V'_{xy'}V_{x'y'}\) с существованием \(xy\)- и \(x'y'\)-вещей и отсутствием \(xy'\)-вещей; вектор \(+-0+\) затем кодирует эти статусы четырёх сочетаний. Поэтому для его модели экзистенциальный смысл не является нашей догадкой. Однако знак \(+\) сам по себе не является квантором \(\exists\): перевод модели в утверждения о конкретных представителях конечной группы — уже применение этой авторской интерпретации к нашему примеру. Ноль не означает, что автору неизвестно, есть ли такие вещи, и не является третьим значением высказывания. \href{https://ternarycomp.cs.msu.ru/Papers/Rea.pdf}{«Реанимация аристотелевой силлогистики» (PDF с. 1–3)}

В примере мы применяем указанную автором экстенсиональную трактовку к конечному списку участников: позиции (+) представлены хотя бы одним человеком, позиция (-) не представлена, а в позиции (0) наличие или отсутствие человека не влияет на профиль. Такое распределение людей — наша конечная иллюстрация; источник прямо обосновывает соответствующий смысл через непустоту классов в принятом универсуме. Н. П. Брусенцов отдельно требует, чтобы для каждого первичного термина в универсуме сосуществовали вещи, которым он присущ, и вещи, которым он антиприсущ; это постулат сосуществования противоположностей.

\begin{longtable}{P{.18\linewidth}P{.43\linewidth}P{.27\linewidth}}
\toprule
\textbf{Сочетание признаков} & \textbf{Пример в собрании} & \textbf{Статус по Н. П. Брусенцову} \\
\midrule
\(xy\): в очках и студент & Анна & \(+\): требуется \\
\(xy'\): в очках и не студент & Таких нет & \(-\): исключается \\
\(x'y\): без очков и студент & Вера & \(0\): присутствие не задаётся \\
\(x'y'\): без очков и не студент & Борис & \(+\): требуется \\
\bottomrule
\end{longtable}

\subsection{Чем профиль отличается от обычного высказывания}

При таком чтении \(xy, xy', x'y, x'y'\) — \textbf{группы предметов}, а не четыре строки таблицы истинности высказываний \(P\) и \(Q\). Профиль \(+-0+\) требует представителей \(xy\) и \(x'y'\), исключает \(xy'\), а наличие \(x'y\) оставляет открытым. Ноль здесь относится к одной из четырёх ячеек, а не к третьему истинностному значению высказывания.

Для записи высказываний введём предикаты: \(X(t)\) — «участник \(t\) носит очки», \(Y(t)\) — «участник \(t\) студент». Здесь переменная \(t\) пробегает только участников этого собрания; \(\forall\) означает «для каждого», \(\exists\) — «существует хотя бы один». В следующей формуле \(\land\) читается «и», а \(\neg\) — «не».

\textbf{Какое обычное высказывание проверяется на собрании?} Фраза «Все участники в очках — студенты» записывается как \(\forall t\,(X(t)\to Y(t))\). Она истинна, если в собрании нет ни одного человека в очках, который не студент, то есть если группа \(xy'\) пуста. Обратная фраза — «Все студенты на собрании носят очки», \(\forall t\,(Y(t)\to X(t))\), — проверяет другое направление. Она ложна в нашем примере из-за Веры: она студентка без очков.

\textbf{Что добавляет критерий Н. П. Брусенцова?} При указанном составе собрания обычная фраза и профиль \(+-0+\), прочитанный в авторском экстенсиональном смысле, оба выполняются. Но это не одна и та же проверка: обычная универсальная фраза запрещает только \(xy'\), тогда как форма отношения в универсуме Н. П. Брусенцова также требует непустоты \(xy\) и \(x'y'\). Если Борис уйдёт, фраза «Все участники в очках — студенты» останется истинной, а профиль перестанет выполняться: в группе нет представителя \(x'y'\). Если же появится не студент в очках, сочетание \(xy'\) будет представлено, и обе проверки не пройдут.

\textbf{Как записать авторское экстенсиональное чтение средствами логики предикатов?} Здесь мы переводим в обычную логику те условия на непустоту и пустоту подклассов, которые Н. П. Брусенцов описывает в главе 2005 года. Знак \(\exists\) означает «существует хотя бы один участник», \(\land\) — «и», \(\neg\) — «не». Формула требует участника в очках и студента, запрещает человека в очках, который не студент, и требует участника без очков, который не студент. Она ничего не требует о студентах без очков.

\[
(\exists t\,(X(t)\land Y(t)))\;\land\; \neg\exists t\,(X(t)\land\neg Y(t))\;\land\; (\exists t\,(\neg X(t)\land\neg Y(t))).
\]

У Н. П. Брусенцова \(x\) и \(y\) — двузначные признаки отдельных предметов: предмет либо носит очки, либо нет; либо студент, либо нет. Приведённая формула — наша запись средствами современной логики предикатов экстенсионального смысла авторской формы для данной области; это не буквальная формула из статьи и не само отношение \(x\Rightarrow y\). Она сама по себе не доказывает, что все операции системы Н. П. Брусенцова эквивалентны классической логике. В статье 2003 года он вводит отдельную функцию \(A(x,y)\) со значениями \(1\), \(0\) и \(\sigma\); её нельзя смешивать ни с привходящим нулём в четырёхъячеечном профиле 2005 года, ни с истинностной оценкой \(P\to Q\).

\paragraph{Граница примера.} Имена участников и их распределение по четырём группам — наша иллюстрация того, как можно прочитать профиль. Она не является формулировкой Н. П. Брусенцова о содержательном следовании и не доказывает его тезис об Аристотеле.

\paragraph{Источник главы 2005 года.} Глава «Реанимация аристотелевой силлогистики» опубликована в сборнике «Реставрация логики» (Москва: фонд «Новое тысячелетие», 2005, с. 140–145). Каталог РГБ указывает для сборника 164 страницы (запись № 002754741). «Неадекватность двоичной информатики» — отдельная статья, где Н. П. Брусенцов ссылается на эту главу. \href{https://www.rsl.ru/photo/%21_ORS/5-PROFESSIONALAM/1-projekt-NEB/spiski_knig_dlya_ozifrovki/2016/doc_8432va.pdf}{РГБ, запись № 1090 в перечне} · \href{https://www.trinitas.ru/rus/doc/0226/002a/02260078.pdf}{Н. П. Брусенцов, «Неадекватность двоичной информатики»}

\subsection{Как не смешивать обозначения 2003, 2005 и 2010 годов}

Похожие обозначения в этих публикациях нельзя считать одной и той же шкалой. В таблице сведены различия, нужные для чтения формул:

\begin{longtable}{P{.28\linewidth}P{.66\linewidth}}
\toprule
\textbf{Запись и источник} & \textbf{Значение} \\
\midrule
\(A(x,y)\in\{1,0,\sigma\}\), статья 2003 года & Знак \(\in\) означает «принадлежит множеству», фигурные скобки перечисляют возможные значения. \(1\) — необходимая присущность, \(0\) — антиприсущность, \(\sigma\) — привходящий статус: \(y\) не необходимо присуще и не антиприсуще \(x\); ни \(xy\)-, ни \(xy'\)-вещи не исключены с необходимостью. Это не просто нехватка знания о фактической группе. Условие \(0<\sigma<1\) задаёт положение \(\sigma\) между числами 0 и 1, но не указывает конкретное значение. \\
\(xy\lor\sigma x'y\lor x'y'\), глава 2005 года & Знак \(\lor\) читается «или» и соединяет выражения о классах; \(\sigma\) помечает привходящую ячейку \(x'y\). В знаковой шкале этой главы \(+1,-1,0\) означают необходимое включение, необходимое исключение и привходящий статус соответственно. Поэтому ноль в профиле \(+-0+\) — не ноль из функции 2003 года, где он означает антиприсущность. \\
\(Axy\) и \(A(x,y)\) & В статье «Воссоздание…» (2010) \(Axy\) означает категорическое суждение «Все \(x\) суть \(y\)». В статье 2003 года \(A(x,y)\) — функция, задающая статус отношения терминов. Это разные объекты, несмотря на похожую запись. \\
\(\text{необходимо / возможно / невозможно}\), статья «Трехзначность отношения следования» (2010) & Это три статуса отношения следования в целом. Их не следует смешивать с тремя статусами отдельных ячеек \(+,-,0\) в профиле 2005 года или со значениями функции \(A(x,y)\) из статьи 2003 года. \\
\(+-++\), наша контрольная схема & Это сравнение точечной материальной импликации \(X(t)\to Y(t)\): здесь \(+\) означает истину, \(-\) — ложь. Эти знаки кодируют истинность формулы, а не статусы Н. П. Брусенцова. \\
\bottomrule
\end{longtable}

\newpage
\subsection{Три статуса в статье 2003 года}

В статье «Трехзначная интерпретация силлогистики Аристотеля» Н. П. Брусенцов называет \(A(x,y)\) двухместной функцией терминов и различает три статуса отношения присущности:

\begin{center}
\begin{tabular}{P{.17\linewidth}P{.34\linewidth}P{.40\linewidth}}
\toprule
\textbf{Значение} & \textbf{Статус у Н. П. Брусенцова} & \textbf{Смысл в модели статьи} \\
\midrule
\(1\) & Необходимо присуще & \(y\) необходимо присущ \(x\); сочетание \(xy'\) исключено. \\
\(0\) & Антиприсуще & Сочетание \(xy\) исключено: возможность присущности \(y\) предметам \(x\) отвергнута. \\
\(\sigma,\;0<\sigma<1\) & Привходящее & \(y\) не необходимо присуще и не антиприсуще \(x\); ни \(xy\)-, ни \(xy'\)-вещи не исключены с необходимостью. Это не просто сообщение о нехватке знания. Неравенство указывает положение символа, но не задаёт конкретное число для \(\sigma\). Смешанная группа с сочетаниями \(xy\) и \(xy'\) — возможный пример, но не определение этого статуса. \\
\bottomrule
\end{tabular}
\end{center}

Функция статьи 2003 года различает три названных \textbf{качественных статуса}; \(\sigma\) — обозначение привходящего случая. Условие \(0<\sigma<1\) помещает его между крайними значениями, но текст не задаёт конкретное число и не позволяет заключить, выбирается ли фиксированное значение или допускаются разные числовые реализации. В главе 2005 года Н. П. Брусенцов также называет привходящий статус третьим значением классового выражения. В обобщённой дизъюнктивной форме \(xy\lor\sigma x'y\lor x'y'\) символ \(\sigma\) помечает ячейку \(x'y\); в знаковой шкале этой главы тот же тип статуса в профиле \(+-0+\) записан нулём. Это соответствие внутри изложения 2005 года не означает равенства \(\sigma=0\) в шкале статьи 2003 года, где \(0\) означает антиприсущность. Ни одна из этих кодировок сама по себе не присваивает третье значение высказыванию \(P\to Q\).

Статья также вводит постулат, который автор называет «сосуществованием противоположностей» или постулатом Гераклита. Для каждого первичного термина он требует сосуществования термина и его противоположности; в двузначном примере «очки / без очков» это означает, что представлены оба случая. В схеме \(V_xV_{x'}V_yV_{y'}V_zV_{z'}\cdots\equiv1\) нижний индекс указывает, какой термин подставлен в выражение \(V\): например, \(V_x\) относится к термину \(x\), а \(V_{x'}\) — к противоположному термину. Пара \(z,z'\) — ещё один пример термина и его противоположности; многоточие означает продолжение цепочки для других терминов. В контексте постулата эти выражения относятся к наличию соответствующих терминов в универсуме. Само \(V\) — специальная операция модели, не обычная дизъюнкция \(\lor\); знак \(1\) — единичное значение авторской алгебры, а не количество предметов. Соседние множители и знак \(\equiv\) читаются в контексте этой алгебры; общей таблицы операций эта формула не задаёт. Сам Н. П. Брусенцов отмечает, что Аристотель явно этот постулат не формулировал. Это добавленная предпосылка реконструкции, а не аксиома Аристотеля. \href{https://computer-museum.ru/histussr/trilog0.htm}{Статья 2003 года в архиве Виртуального компьютерного музея}

Обозначения и операции в формуле необходимого следования из статьи 2010 года, а также различие между условием универсума и профилем сочетаний, разобраны отдельно в руководстве \href{Как-читать-формулы-Брусенцова.tex}{«Как читать формулы Н. П. Брусенцова»}.

\subsection{Закон исключённого третьего и «средний» случай}

В статье 2003 года Н. П. Брусенцов пишет, что закон «несправедливо приписан» Аристотелю как хрисиппов постулат (PDF с. 1; точные слова автора: «несправедливо приписав ему хрисиппов постулат исключенного третьего»). Это тезис автора статьи, а не доказательство истории эллинистической логики. В «Метафизике» Γ.7, 1011b23–25, в буквальном переводе W. D. Ross сказано: «между противоречащими не может быть промежуточного; об одном предмете мы должны либо утверждать, либо отрицать любой предикат» (английский текст и греческий оригинал доступны в указанном читателе). Это прямое основание считать, что Аристотель формулирует принцип, позднее записываемый как \(P\lor\neg P\); оно само по себе не доказывает, что он вводил закон как отдельную формальную аксиому или употреблял позднейшее название. \href{https://aristotle.johnhboyer.com/Meta/book/4/}{Аристотель, «Метафизика», Γ.7, 1011b23–25; перевод W. D. Ross}

При этом закон исключённого третьего и принцип двузначности — связанные, но разные формулировки: первый утверждает \(P\lor\neg P\), второй требует, чтобы каждое высказывание имело ровно одно из двух значений — истину или ложь. В «Об истолковании» 9 Аристотель отдельно обсуждает будущие случайные события: необходимо, что морское сражение завтра будет или не будет, но ни один из этих исходов не является необходимым сам по себе. Это ограничивает простой перенос двузначной оценки на будущие высказывания, однако не доказывает, что статус \(\sigma\) у Н. П. Брусенцова — третье истинностное значение Аристотеля. \href{https://www.st-andrews.ac.uk/~mnat/~ball0888/salamis/interpretatione.html}{Аристотель, «Об истолковании», гл. 9}

Случай \(\sigma\) можно понять без нарушения закона исключённого третьего. Смешанная группа, где часть предметов \(x\) обладает признаком \(y\), а часть — нет, служит одним из примеров, когда не выполнено ни отношение «всякое \(x\) есть \(y\)», ни отношение «ни одно \(x\) не есть \(y\)». Но такая смешанность — иллюстрация, а не само определение \(\sigma\) у Н. П. Брусенцова. В статье 2003 года привходящий статус определяется через отсутствие необходимого включения и необходимого исключения. Поэтому \(\sigma\) описывает промежуточный статус отношения терминов, а не промежуток между высказыванием \(P\) и его отрицанием \(\neg P\).

Н. П. Брусенцов также заявляет, что двухзначная логика не позволяет совместить подчинение частного суждения общему с законом исключённого третьего. Но сама двузначность этого не запрещает: при явно заданном существовании предметов класса \(X\) из общего суждения получается частное:

\[
\forall t\,(X(t)\to Y(t)),\quad \exists t\,X(t)
\;\vdash\; \exists t\,(X(t)\land Y(t)).
\]

Непустота класса \(X\) здесь выступает отдельной посылкой (экзистенциальным импортом). Поэтому спор касается условий существования терминов и выбранной формализации, а не логической несовместимости подчинения с законом исключённого третьего.

\subsection{Что подтверждает PDF-копия}

Формулы в основном тексте статьи читаются на изображениях PDF, но копия была пересобрана из архивного Word-файла и не является факсимиле. В сплошном тексте местами пропадают штрихи у \(x'\) и \(y'\); в текстовом извлечении штрих, \(\lor\), \(\land\), \(\equiv\) и \(\sigma\) могут превращаться в посторонние символы. Страница с диаграммами отображается с повреждёнными подписями и наложением элементов, поэтому по этой копии нельзя надёжно проверить все обозначения на рисунках.

Есть и расхождение внутри статьи: русский основной текст сообщает 192 правильных модуса, тогда как английская аннотация указывает 190. В других работах Н. П. Брусенцова встречается число 128. Пока основания подсчёта не сопоставлены, ни одно из чисел не следует представлять как бесспорное.

\section{Когда связь можно назвать содержательной}

Для разбора математического доказательства полезен вопрос: \textbf{какое сведение даёт посылка и как оно обосновывает заключение?} Из условия «\(n\) делится на 4» мы получаем целый множитель \(k\), для которого \(n=4k\). Через него находим целое число \(2k\), удвоение которого даёт \(n\). Именно это требовалось для вывода о чётности. Совпадения темы «числа» для такого вывода мало: из посылки «\(n>10\)» чётность \(n\) не следует. Контрпример — число 11.

\begin{longtable}{P{.23\linewidth}P{.70\linewidth}}
\toprule
\textbf{Если нас интересует…} & \textbf{Какой вопрос решается} \\
\midrule
\endfirsthead
\toprule
\textbf{Если нас интересует…} & \textbf{Какой вопрос решается} \\
\midrule
\endhead
Классическое следование & Есть ли допустимая модель, в которой все посылки истинны, а заключение ложно? Одной такой модели достаточно, чтобы опровергнуть следование. \\
Строгая импликация в модальной логике & Утверждение «необходимо, что если P, то Q» истинно в данном мире, когда среди доступных из него миров нет случая, где P истинно, а Q ложно. Например, невозможность P уже исключает такой случай для любого Q. \\
Релевантное следование & Как в данной системе обеспечивается участие посылки в выводе? Релевантные логики исключают, например, ряд схем, связывающих противоречие с произвольным посторонним заключением. \\
Объяснение факта & Какое определение, свойство или причина объясняет, почему заключение верно о данном предмете? \\
\bottomrule
\end{longtable}

В примере с числами определение делимости объясняет, почему контрпример невозможен. В примере с собранием достаточно проверить список, чтобы установить, что контрпример отсутствует \textbf{у данных участников}. Возможное появление человека в очках, который не является студентом, показывает предел такого обобщения. \href{https://plato.stanford.edu/entries/logic-conditionals/}{Исследования условных высказываний} · \href{https://plato.stanford.edu/entries/logic-relevance/}{Релевантная логика}

\section{Короткая карта обозначений}

\begin{longtable}{P{.18\linewidth}P{.43\linewidth}P{.27\linewidth}}
\toprule
\textbf{Запись} & \textbf{Что она обозначает} & \textbf{Какой вопрос задаёт} \\
\midrule
\(P\to Q\) & Условное высказывание («если \(P\), то \(Q\)») & Истинно ли это высказывание при заданных значениях \(P\) и \(Q\)? \\
\(P\models Q\) & Семантическое следование & Во всех ли моделях, где посылка \(P\) истинна, истинно и \(Q\)? \\
\(P\vdash Q\) & Выводимость в исчислении & Можно ли получить \(Q\) из посылок по правилам этого исчисления? \\
\(P\to Q,\ P\vdash Q\) & Применение \emph{modus ponens} & Разрешает ли правило из условной посылки и \(P\) записать заключение \(Q\)? Да. \\
\(x\Rightarrow y\) у Н. П. Брусенцова & Отношение необходимого содержательного следования \(y\) из \(x\); профиль \(+-0+\) — одно из его алгебраических представлений & Как автор формализует необходимое следование между терминами и какой статус имеет каждое сочетание \(x/y\)? \\
\bottomrule
\end{longtable}

Заключение \(Q\) — это высказывание, которое получено. Оно не является ни \(P\to Q\), ни правилом вывода. Импликация может быть посылкой или способом выразить условие; правило объясняет сам разрешённый шаг.

\newpage
\section{Источники}

\begin{itemize}
\item Н. П. Брусенцов. \href{https://computer-museum.ru/histussr/trilog0.htm}{«Трехзначная интерпретация силлогистики Аристотеля»} (2003). Архивная публикация; локальная PDF-копия пересобрана из DOC и не является факсимиле.
\item Н. П. Брусенцов. \href{https://ternarycomp.cs.msu.ru/Papers/Rea.pdf}{«Реанимация аристотелевой силлогистики»} (2005), глава из книги «Реставрация логики», с. 140–145.
\item Н. П. Брусенцов. \href{https://ternarycomp.cs.msu.ru/Papers/Soder_sled.pdf}{«Содержательное следование, сосуществование противоположностей и формальные импликации»} (2006), с. 342–343.
\item Н. П. Брусенцов. \href{https://ternarycomp.cs.msu.ru/Papers/Alrek.pdf}{«Алгебраическая реконструкция силлогистики»} (2008), с. 7–10.
\item Н. П. Брусенцов. \href{https://ternarycomp.cs.msu.ru/Papers/brusentsov_sborn2010_1.pdf}{«Трехзначность отношения следования»} (2010), с. 86–87.
\item Н. П. Брусенцов. \href{https://ternarycomp.cs.msu.ru/Papers/brusentsov_sborn2010_2.pdf}{«Воссоздание аристотелевой безукоризненной силлогистики»} (2010), с. 88–91.
\item Аристотель. \href{https://classics.mit.edu/Aristotle/prior.1.i.html}{«Первая аналитика»}, I.1.
\item Аристотель. \href{https://classics.mit.edu/Aristotle/posterior.1.i.html}{«Вторая аналитика»}, I.2.
\item Аристотель. \href{https://aristotle.johnhboyer.com/Meta/book/4/}{«Метафизика»}, Γ.7.
\item Аристотель. \href{https://www.st-andrews.ac.uk/~mnat/~ball0888/salamis/interpretatione.html}{«Об истолковании»}, гл. 9.
\item George Boole. \href{https://www.gutenberg.org/ebooks/36884}{\emph{The Mathematical Analysis of Logic} (1847)}.
\item George Boole. \href{https://www.gutenberg.org/ebooks/15114}{\emph{An Investigation of the Laws of Thought} (1854)}.
\item Lewis Carroll. \href{https://fair-use.org/mind/1895/04/what-the-tortoise-said-to-achilles}{“What the Tortoise Said to Achilles”}, \emph{Mind} 4, no. 14 (1895), 278–280.
\item Lewis Carroll. \href{https://archive.org/details/symboliclogicpa00carrgoog}{\emph{Symbolic Logic, Part I: Elementary} (1896)}.
\item Internet Encyclopedia of Philosophy. \href{https://iep.utm.edu/lewis-carroll-logic/}{Lewis Carroll: Logic}.
\item Stanford Encyclopedia of Philosophy. \href{https://plato.stanford.edu/entries/algebra-logic-tradition/}{The Algebra of Logic Tradition}.
\item Stanford Encyclopedia of Philosophy. \href{https://plato.stanford.edu/entries/logical-consequence/}{Logical Consequence}, \href{https://plato.stanford.edu/entries/logic-conditionals/}{Logic and Conditionals}, \href{https://plato.stanford.edu/entries/logic-relevance/}{Relevance Logic}, \href{https://plato.stanford.edu/entries/aristotle-logic/}{Aristotle’s Logic}.
\end{itemize}

\end{document}
